\paragraph{\trIdempotent{}元与 A∞-模。}\label{sub:idempotents-modules}
把基域 $\KK$ 以显然的方式看成一个只有一个\trObject{} $*$ 的 \trACategory。当然，
它的\trCohomology\trCategory{}又是 $\KK$，不过现在把它看成通常的\trCategory{}罢了。
我们不会在记号上区分这两者；这应当只会造成极小的混淆。

设 $\scr{A}$ 是一个 \trACategory，$A=H(\scr{A})$ 是它的\trCohomology\trCategory，
而 $\scr{Q}=\mathit{mod}(\scr{A})$ 是 \trAinfty-模的\trCategory。一个\trObject{} $Y$
的\textbf{\trHomotopy\trIdempotent}（idempotent up to homotopy）定义为一个%
\trNonunital{} \trAFunctor{} $\wp\colon\KK\to\scr{A}$，满足 $\wp(*)=Y$。具体地，它由
元素 $\wp^d\in\hom^{1-d}_{\scr{A}}(Y,Y)$（$d\ge 1$）给出，这些元素满足方程组
\begin{equation}\label{eq:homotopy-idempotent-equations}
\sum_r\sum_{s_1,\ldots,s_r}
\mu^r_{\scr{A}}\bigl(\wp^{s_1},\ldots,\wp^{s_r}\bigr)
=
\begin{cases}
\wp^{d-1} & \text{若 } d \text{ 为偶数，}\\
0 & \text{若 } d \text{ 为奇数}
\end{cases}
\end{equation}
其中求和取遍\trPartition{} $s_1+\cdots+s_r=d$。

\begin{lemma}\label{thm:homotopy-idempotent-lift}
设 $p\in\Hom_A(Y,Y)$ 是\trCohomology\trCategory{}中的一个\trIdempotent{}%
\trEndomorphism。那么存在一个\trHomotopy\trIdempotent{} $\wp$，使得
$[\wp^1]=p$。
\end{lemma}
\begin{proof}
$p$ 定义了一个\trNonunital{} \trFunctor{} $\KK\to A$。据%
\cref{thm:functor-realization-obstruction}，把它\trLift{}成一个 \trAFunctor{}的%
\trObstruction{}落在 $HH^2(\KK,A)^{2-r}$ 中（$r>2$）。其底层的%
\trHochschild{}\trComplex{}是
\begin{equation*}
CC^{r+s}(\KK,A)^s=\Hom^s_A(Y,Y),\qquad
\partial^{r,s}(T)=(-1)^s
\begin{cases}
pT-Tp & \text{若 } r \text{ 为偶数，}\\
pT-Tp^\perp & \text{若 } r \text{ 为奇数。}
\end{cases}
\end{equation*}
把 $\Hom_A(Y,Y)$ 按与 $p$ 的左、右\trComposition{}的\trImage{}与\trKernel{}%
拆分之后，立刻可以看出仅有的非零\trCohomology{}群是
$HH^s(\KK,A)^s=p\,\Hom^s_A(Y,Y)\,p\oplus p^\perp\,\Hom^s_A(Y,Y)\,p^\perp$。
\end{proof}

\begin{remark}\label{thm:homotopy-idempotent-ext}
与其用 \trHochschild{}\trComplex{}作直接计算，也可以用一些基本的%
\trHomology{}代数。先给 $\KK$
添加一个幺元，使得所讨论的\trFunctor{}变成\trUnital{}的；如此所得的、如%
\cref{thm:morita-bimod-route} 中的\trBimodule{}\trCategory{}是\trSemisimple{}的，
因此高阶 $\operatorname{Ext}$ 群全部消没。
\end{remark}

设 $\wp$ 是一个\trHomotopy\trIdempotent。我们给它配上一个 \trAinfty-模
$\scr{Z}$，称为 $\wp$ 的\textbf{抽象像}（abstract image），其定义如下：
\begin{equation*}
\scr{Z}(X)=\hom_{\scr{A}}(X,Y)[q]
\end{equation*}
是由单个\trGraded{}变量 $q$（$|q|=-1$）的形式多项式组成的空间，其%
\trCoefficients{}取自 $\hom_{\scr{A}}(X,Y)$。记 $\delta_q$ 为归一化的形式求导
运算，它把 $q^0$ 送到 $0$，而把其余各次幂 $q^k$ 送到 $q^{k-1}$。采用这一记号，
\trDifferential{}便是
\begin{equation}\label{eq:abstract-image-differential}
\begin{split}
\mu^1_{\scr{Z}}(b(q))
&=\sum_r\sum_{s_2,\ldots,s_r}\delta_q^{s_2+\cdots+s_r}\,
\mu^r_{\scr{A}}\bigl(\wp^{s_r},\ldots,\wp^{s_2},b(q)\bigr)\\
&\quad+(-1)^{|b|}\frac{b(q)-b(-q)}{2q},
\end{split}
\end{equation}
其中最后一项是先把幂级数\trAntisymmetric{}化、再施以 $\delta_q$；这并不真正
涉及乘以 $\frac{1}{2}$，因此当 $\KK$ 的特征为 $2$ 时也有意义。高阶%
\trStructureMap{}（对 $d\ge 2$）为
\begin{equation*}
\begin{split}
&\mu^d_{\scr{Z}}\bigl(b(q),a_{d-1},\ldots,a_1\bigr)\\
&\qquad=\sum_r\sum_{s_{d+1},\ldots,s_r}\delta_q^{s_{d+1}+\cdots+s_r}\,
\mu^r_{\scr{A}}\bigl(\wp^{s_r},\ldots,\wp^{s_{d+1}},b(q),a_{d-1},\ldots,a_1\bigr).
\end{split}
\end{equation*}

\begin{lemma}\label{thm:abstract-image-homology}
底层\trCohomology{}模 $H(\scr{Z})$ 可以描述如下：$H(\scr{Z})(X)\cong
p\,\Hom_A(X,Y)$，其中 $p=[\wp^1]\in\Hom_A(Y,Y)$，而且\trModule{}结构
$H(\scr{Z})(X_1)\otimes\Hom_A(X_0,X_1)\to H(\scr{Z})(X_0)$%
\footnote{译注：原书此处印作 $\Hom(\scr{Z})(X_0)$；疑应为 $H(\scr{Z})(X_0)$
（同句所述者为底层\trCohomology{}模 $H(\scr{Z})$ 的\trModule{}结构，且该映射
左端即为 $H(\scr{Z})(X_1)$）；此处径改。——译者注}来自通常的\trComposition{}%
$p\,\Hom_A(X_1,Y)\otimes\Hom_A(X_0,X_1)\to p\,\Hom_A(X_0,Y)$。
\end{lemma}
\begin{proof}
\begingroup\emergencystretch=8em\hbadness=10000
$\scr{Z}(X)$ 有一个自然的递增\trFiltration，即
$F_r=\hom_{\scr{A}}(X,Y)q^0\oplus\cdots\oplus$
$\hom_{\scr{A}}(X,Y)q^{r-1}$。我们
\par\endgroup
通常用的是递减\trFiltration，为统一口径起见，令 $r\leftrightarrow-r$ 把它变成
一个递减的。于是所得的\trSpectralSequence{}具有
\begin{equation*}
E_1^{-r,s}=H^{s-r}(F_{r+1}/F_r)=
\begin{cases}
\Hom^s_A(X,Y), & r\ge 0,\\
0, & \text{其余情形。}
\end{cases}
\end{equation*}
其\trDifferential{}由\eqref{eq:abstract-image-differential} 中恰含一个 $\delta_q$
或 $1/q$ 的那些项所诱导，故由下式给出：
\begin{equation*}
\partial^{-r,s}(b)=(-1)^{r+s}
\begin{cases}
p^\perp b, & r>0 \text{ 为奇数},\\
pb, & r>0 \text{ 为偶数},\\
0, & \text{其余情形。}
\end{cases}
\end{equation*}
于是，该\trSpectralSequence{}在 $E_2$ 处即已\trDegenerate，而 $H(\scr{Z})(X)$
\trIsomorphic{}于唯一非零的列
$E_2^{0\bullet}=p\,\Hom_A(X,Y)$。事实上，此\trSpectralSequence{}的%
边缘\trHomomorphism{}给出这两个群之间的一个\trCanonical{}\trIsomorphism；
据此即可看出，与\trModule{}结构所需的\trCompatible{}性是显然的。
\end{proof}

\begin{lemma}\label{thm:hom-abstract-image}
对任何 \trAinfty-模 $\scr{M}$，存在\trIsomorphism
\begin{equation*}
\Hom_{H(\scr{Q})}(\scr{Z},\scr{M})\cong H(\scr{M})(Y)\,p,
\end{equation*}
而且它关于用\trModule\trHomomorphism{}作左\trComposition{}是自然的。
\end{lemma}
\begin{proof}
考虑\cref{thm:yoneda-ext-interpretation}中的阿贝尔\trCategory{}%
$\mathit{Mod}^0(A)$。\cref{thm:abstract-image-homology}说，作为该\trCategory{}%
的一个\trObject，$H(\scr{Z})$ 是 $\Hom_A(-,Y)$ 的一个\trDirectSummand；而既然
后者是投射的，前者亦然。鉴于\eqref{eq:ext-e2-term}，这意味着：由
$\hom_{\scr{Q}}(\scr{Z},\scr{M})$ 上的长度\trFiltration{}所诱导的%
\trSpectralSequence{}在 $E_2$ 处\trDegenerate，而且映射
$\Hom_{H(\scr{Q})}(\scr{Z},\scr{M})\to
\Hom_{\mathit{Mod}^0(A)}\bigl(H(\scr{Z}),H(\scr{M})[\bullet]\bigr)$
是一个\trIsomorphism。由经典的\trYoneda{}引理连同%
\cref{thm:abstract-image-homology}，后一个\trGraded{}群自然地\trIsomorphic{}于
$H(\scr{M})(Y)\,p$。
\end{proof}

考虑对 $p$ 施以\trYonedaEmbedding{}而得到的 $H^0(\scr{Q})$ 中的%
\trIdempotent{}\trEndomorphism。\cref{thm:hom-abstract-image}直接蕴含：$\scr{Z}$
就是这个\trEndomorphism{}在\eqref{eq:image-inclusion-retraction} 意义下的%
\trImage。特别地，在相差 $H^0(\scr{Q})$ 中一个\trIsomorphism{}的意义下，抽象像
只依赖于 $p$，而不依赖于 $\wp$ 的选取。

\begin{remark}\label{thm:abstract-image-history}
$\scr{Z}$ 的定义是受 [24] 中的一个观察启发而来的；该观察说，在合适的%
\trComplex{}\trCategory{}中，\trDirectSummand{}可以凭一个无穷分解技巧而存在
（感谢 Fran\c{c}ois Loeser 向我指出这一点）。
\end{remark}
